%% ==========================================================================
%%  Approach 12 --- a complete, self-contained proof of Conjecture 2 at n = 3.
%%
%%  Source: RESIDUAL.md \S7.16.32--7.16.33; formalised piecemeal in
%%  \Cref{sec:approach11}. This section restates the same result as one
%%  linear proof, written to be read on its own: every definition it needs is
%%  given here, and every lemma kept is one the final proof actually calls.
%%  Nothing about paths that were tried and abandoned appears below; see
%%  RESIDUAL.md \S7 and \Cref{sec:approach10} for that history.
%% ==========================================================================
\section{Approach 12: A Complete Proof of Conjecture 2 at \texorpdfstring{$n=3$}{n=3}}
\label{sec:approach12}

This section proves Conjecture~\ref{conj:main} for $n=3$ agents, in full
generality, from scratch. It assumes nothing from the rest of this document
except the two definitions restated in \Cref{sec:g2-setup} and standard
facts about set functions. It is self-contained on purpose: everything a
reader needs is either defined below or is a short, complete proof.

\subsection{Setup}
\label{sec:g2-setup}

There are $3$ agents and a finite set $\items$ of indivisible chores. Each
agent $i$ has a \emph{dichotomous cost function} $\cost_i : 2^{\items} \to
\R_{\ge0}$: $\cost_i(\emptyset) = 0$, and every marginal $\cost_i(S \cup
\set e) - \cost_i(S)$, for $S \subseteq \items$ and $e \notin S$, lies in
$\set{0,1}$. This forces $\cost_i$ to be monotone ($\cost_i(S) \le
\cost_i(T)$ whenever $S \subseteq T$), but assumes nothing else: not
additivity, not submodularity, not any relationship between $\cost_1,
\cost_2, \cost_3$.

An \emph{allocation} $\alloc = (\bundle1,\bundle2,\bundle3)$ is a partition
of $\items$ into three (possibly empty) bundles, one per agent. A subsidy
vector $\subsidy = (\subsidy_1,\subsidy_2,\subsidy_3) \in \set{0,1}^3$ pays
agent $i$ the amount $\subsidy_i$. The pair $(\alloc,\subsidy)$ is
\emph{envy-free} (EF) if
\[
  \cost_i(\bundle i) - \subsidy_i \;\le\; \cost_i(\bundle j) - \subsidy_j
  \qquad \text{for all agents } i,j:
\]
net of compensation, no agent would rather have someone else's bundle and
subsidy. A \emph{partial allocation} assigns each item to at most one agent
(some items may be unallocated); it is EF if the inequality above holds for
all $i,j$ using only the bundles actually assigned so far.

\begin{quote}
\textbf{Conjecture 2.} For every instance of $3$ dichotomous cost functions
on a finite item set $\items$, there is an allocation $\alloc$ and a subsidy
$\subsidy \in \set{0,1}^3$ such that $(\alloc,\subsidy)$ is envy-free, and
such a pair can be found in polynomial time.
\end{quote}

The rest of this section proves this in five steps. Read top to bottom; each
step uses only what came before it.

\subsection{Step 1: start from a partial allocation that is already free}
\label{sec:g2-step1}

The proof does not build an allocation from nothing. It starts from an
external result that hands us most of one for free.

\begin{theorem}[Tao, Wu, Yu, Zhou 2023~\cite{TWYZ23}]
\label{thm:g2-twyz}
For any number of agents $n$ and any dichotomous cost functions, there is a
polynomial-time algorithm producing a partial allocation
$(\bundle1,\dots,\bundle n)$ that is envy-free with $\subsidy = 0$ (no
subsidy needed at all) and leaves at most $n-1$ items of $\items$
unassigned.
\end{theorem}

We use this as a black box; the idea behind it (an envy-graph process that
gives an item to whoever is willing to take it for free, and otherwise
rotates bundles around a cycle of mutually-envious agents) is not needed
here. At $n=3$, \Cref{thm:g2-twyz} gives an EF partial allocation
$(\bundle1,\bundle2,\bundle3)$ together with a leftover set $R = \items
\setminus (\bundle1 \cup \bundle2 \cup \bundle3)$ of at most $2$ items. If
$R = \emptyset$ we are already done, with $\subsidy = 0$. What remains is
to show that the leftover $1$ or $2$ items can always be placed --- possibly
by reassigning which agent gets which of the three bundles, and possibly
with a subsidy of $1$ for one or two agents --- so that the full allocation
is envy-free.

\subsection{Step 2: free items can always be inserted}
\label{sec:g2-step2}

\begin{lemma}[Free insertion]
\label{lem:g2-free}
Let $(\bundle1,\bundle2,\bundle3)$ be an EF partial allocation and let $S$ be
a set of unallocated items with $\cost_i(\bundle i \cup S) = \cost_i(\bundle
i)$ for some agent $i$ (item(s) $S$ cost $i$ nothing on top of what $i$
already holds). Giving all of $S$ to $i$, and changing nothing else, keeps
the allocation EF, with $\subsidy = 0$.
\end{lemma}
\begin{proof}
Agent $i$'s own cost is unchanged, so $i$'s non-envy of every other agent
$j$ is unchanged: $\cost_i(\bundle i \cup S) = \cost_i(\bundle i) \le
\cost_i(\bundle j)$. For any other agent $j$: before the change,
$\cost_j(\bundle j) \le \cost_j(\bundle i)$ (that's $j$ not envying $i$,
which held already); after the change $j$ must compare against $\bundle i
\cup S$ instead of $\bundle i$, but $\cost_j(\bundle i \cup S) \ge
\cost_j(\bundle i)$ by monotonicity, so $\cost_j(\bundle j) \le \cost_j(\bundle
i) \le \cost_j(\bundle i \cup S)$ still holds. No agent's envy status
changes for the worse.
\end{proof}

This lemma costs nothing to apply and needs no case analysis, so the proof
always uses it first: whenever some leftover item is free to some agent's
current bundle, hand it over and shrink the leftover set by one. Repeating
this can only ever reduce $\abs R$. The only way this process does not
immediately finish the job is if it gets stuck: every agent has a strictly
positive marginal cost for every leftover item, on their own current bundle.
Call this the \emph{stuck} case. It is the only case left to handle, and it
splits into $\abs R = 1$ and $\abs R = 2$.

\subsection{A motivating example: subsidy alone is not enough}
\label{sec:g2-example}

Before handling the stuck case in general, one example shows why a natural
simpler idea --- keep everyone's bundle as is, place the leftover item, and
patch any envy with a subsidy --- can fail outright, and what actually fixes
it.

\begin{example}
\label{ex:g2-motivating}
Let $\items = \set{a,b,c,e}$ and $\bundle1 = \set a$, $\bundle2 = \set b$,
$\bundle3 = \set c$ (so $R = \set e$). Define, for each agent $i \in
\set{1,2,3}$ with trigger pair $T_1 = \set{a,e}$, $T_2 = \set{b,e}$, $T_3 =
\set{c,e}$,
\[
  \cost_i(S) = \begin{cases} 1 & \text{if } T_i \subseteq S \\ 0 & \text{otherwise.} \end{cases}
\]
Each $\cost_i$ is dichotomous (it only ever jumps by $1$, exactly when the
last missing element of $T_i$ is added). The starting partial allocation is
EF: every agent's cost for every one of the three bundles is $0$, since none
contains any full trigger pair. Every agent has marginal cost $1$ for $e$ on
their own bundle ($\cost_i(\bundle i \cup e) = 1 \ne 0$), so \Cref{lem:g2-free}
does not apply: this is a stuck instance.

Try giving $e$ to agent $1$, keeping bundles otherwise fixed:
$\bundle1' = \set{a,e}$. Now $\cost_1(\bundle1') = 1 > 0 = \cost_1(\bundle2)$,
so agent $1$ envies agent $2$ unless $\subsidy_1 - \subsidy_2 \ge 1$, forcing
$\subsidy_1 = 1, \subsidy_2 = 0$. But agent $2$ sees $\cost_2(\bundle1') = 0$
(agent $2$'s trigger is $\set{b,e}$, and $b \notin \bundle1'$) and
$\cost_2(\bundle2) = 0$: with $\subsidy_1 = 1$, agent $2$'s comparison becomes
$0 - 0 \le 0 - 1$, which is false. Agent $2$ now envies agent $1$. Trying the
other two agents as the recipient of $e$ fails the same way by symmetry, and
no subsidy vector in $\set{0,1}^3$ rescues any of the three choices.

What does work is reassigning bundles, not just placing $e$: give agent $3$
the bundle $\set{a,e}$, agent $1$ the bundle $\set b$, and agent $2$ the
bundle $\set c$. Every agent's cost for their new bundle is $0$ (agent $3$'s
trigger is $\set{c,e}$, and $\set{a,e}$ doesn't contain $c$; similarly for
the other two), so $\subsidy = 0$ works.
\end{example}

The move that rescued \Cref{ex:g2-motivating} was not moving an item between
bundles --- it was reassigning which \emph{agent} receives which
\emph{bundle}. This is the extra freedom the rest of the proof uses: after
deciding where the leftover item(s) go, we are also free to permute which of
the three agents receives which of the three resulting bundles, before
falling back on subsidy.

\subsection{Step 3: a finite question, and why a bounded search decides it}
\label{sec:g2-saturation}

Both remaining cases ($\abs R = 1$ and the stuck case at $\abs R = 2$)
reduce to the same kind of question: given a $3\times3$ table of costs
(agent $i$'s cost for candidate bundle $j$), does some permutation of
bundles to agents plus a subsidy in $\set{0,1}^3$ make the allocation EF?
The next observation turns this into something a finite check can decide
completely, for every possible cost table at once.

\begin{lemma}[Saturation]
\label{lem:g2-saturation}
Let $D$ be a $3\times3$ table, $D_{ij}$ the cost to agent $i$ of candidate
bundle $j$. Whether some permutation $\pi$ of bundles to agents and some
$\subsidy \in \set{0,1}^3$ make the allocation ``agent $i$ gets bundle
$\pi(i)$'' envy-free depends on $D$ only through, for every row $i$ and pair
of columns $j,k$: whether $D_{ij} - D_{ik}$ is $\le -1$, exactly $0$,
exactly $1$, or $\ge 2$. Consequently, if the diagonal of a base matrix $C$
is $0$ (always achievable, see the proof) and its off-diagonal entries are
confined to $\set{0,1,2}$, checking every such $C$ exhaustively already
covers every case an unbounded cost table could ever produce.
\end{lemma}
\begin{proof}
Agent $i$ (holding column $j$) does not envy column $k$ exactly when
$D_{ij} - \subsidy_j \le D_{ik} - \subsidy_k$, i.e.\ $D_{ij} - D_{ik} \le
\subsidy_j - \subsidy_k$. Since $\subsidy \in \set{0,1}^3$, the quantity
$\subsidy_j - \subsidy_k$ only ever takes the values $-1,0,1$. So: if
$D_{ij}-D_{ik} \le -1$ the inequality holds no matter what $\subsidy$ is;
if $=0$ it holds iff $\subsidy_j \ge \subsidy_k$; if $=1$ it holds iff
$\subsidy_j = 1,\subsidy_k = 0$ exactly; if $\ge2$ it never holds. Those are
exactly the four cases claimed.

For the second part: given any cost table, subtract $C_{ii}$ from every
entry in row $i$ (each row independently). This does not change any
difference $D_{ij}-D_{ik}$ within a row, hence does not change which of the
four cases above we are in for any comparison, so it does not change the
answer to the yes/no question --- it only relabels the diagonal to $0$. With
$C_{ii}=0$, the off-diagonal entries are $\ge0$ automatically (they were
$\ge C_{ii}$ before the shift, since agent $i$ never has negative cost
relative to their own bundle in the partial allocations this proof
considers). Restricting them to $\set{0,1,2}$ still lets differences
$C_{ij}-C_{ik}$ range over $\set{-2,\dots,2}$, which covers all four cases
above ($\le-1$ via $-1$ or $-2$, $=0$, $=1$, $\ge2$ via $+2$). An exhaustive
search over $\set{0,1,2}$ therefore realises every combination of the four
cases, across every row and pair, that any cost table --- bounded or not ---
could ever exhibit.
\end{proof}

In short: subsidies can only ever bridge a gap of exactly $1$, so nothing
about a gap larger than $1$ matters beyond the fact that it \emph{is} larger
than $1$, and nothing about the exact size of a gap matters at all once it
exceeds $1$. This is the only idea needed to turn the next two steps from an
open-ended search into two finite, complete checks.

\subsection{Step 4: closing the case of one leftover item}
\label{sec:g2-r1}

\begin{theorem}
\label{thm:g2-r1}
If $(\bundle1,\bundle2,\bundle3)$ is an EF partial allocation with exactly
one leftover item $e$, some choice of which bundle receives $e$, some
permutation of the resulting three bundles to the three agents, and some
$\subsidy \in \set{0,1}^3$ gives a full EF allocation.
\end{theorem}
\begin{proof}
If \Cref{lem:g2-free} applies (some agent has $0$ marginal cost for $e$),
it settles the claim with $\subsidy=0$. Otherwise every agent has marginal
cost exactly $1$ for $e$ on their own bundle. By \Cref{lem:g2-saturation},
after the row shift the entire yes/no question is determined by: the $6$
base off-diagonal entries $C_{ij} \in \set{0,1,2}$ ($\cost_i(\bundle{j})$,
shifted), and $6$ marginal bits $\mu_{ij} = \cost_i(e \mid \bundle{j}) \in
\set{0,1}$ for $i \ne j$ (agent $i$'s marginal cost for $e$ on agent $j$'s
bundle; $\mu_{ii}=1$ is forced by the case we are in). There are $3^6 \times
2^6 = 46{,}656$ such combinations. For each one, checking all $3$ choices of
which bundle absorbs $e$, all $6$ permutations of $3$ bundles to $3$ agents,
and all $8$ subsidy vectors in $\set{0,1}^3$ finds an EF allocation in every
single one --- zero exceptions. (Script: \texttt{update\_12/prove\_r1.py};
runtime under a second. Cross-checked with the off-diagonal range widened to
$\set{0,\dots,3}$ and $\set{0,\dots,4}$ --- $262{,}144$ and $1{,}000{,}000$
combinations respectively, same result, exactly as \Cref{lem:g2-saturation}
predicts.)
\end{proof}

\subsection{Step 5: closing the case of two leftover items}
\label{sec:g2-r2}

\begin{theorem}
\label{thm:g2-r2}
If $(\bundle1,\bundle2,\bundle3)$ is an EF partial allocation with exactly
two leftover items, some assignment of them to bundles, some permutation of
the resulting three bundles to the three agents, and some $\subsidy \in
\set{0,1}^3$ gives a full EF allocation.
\end{theorem}
\begin{proof}
Write $R = \set{e_1,e_2}$. If \Cref{lem:g2-free} applies to either item on
its own, hand it over and the leftover set shrinks to one item;
\Cref{thm:g2-r1} then finishes the job (it does not care how the partial
allocation it receives was produced). So the only case left is where both
items are stuck: every agent has marginal cost $1$ for \emph{both} $e_1$ and
$e_2$ on their own original bundle.

Send $e_1$ to one bundle and $e_2$ to a \emph{different} bundle (never both
to the same bundle --- this already turns out to be enough). By
\Cref{lem:g2-saturation}, the question is again decided by finitely many
bounded quantities: the $6$ base entries $C_{ij} \in \set{0,1,2}$, and now
$12$ marginal bits $\mu_{i,e,j} = \cost_i(e \mid \bundle{j})$ for $i \ne j$,
$e \in \set{1,2}$ (with $\mu_{i,e,i}=1$ forced). There are $3^6 \times
2^{12} = 2{,}985{,}984$ combinations. For each, checking all $6$ ordered
choices of which two (distinct) bundles receive $e_1$ and $e_2$, all $6$
permutations, and all $8$ subsidy vectors finds an EF allocation every time
--- zero exceptions (\texttt{update\_12/prove\_r2.py}, runtime about $20$
seconds).

As an independent check, \texttt{update\_12/verify\_real.py} generates
genuine random dichotomous cost triples (not the bounded abstract model),
finds real EF partial allocations with two leftover items where both are
genuinely stuck for every agent (this is common, not a corner case: with the
script's defaults, over $9{,}000$ such instances turned up in $300$ random
trials), and confirms the same rescue works on every one of them --- and,
in the same run, finds and rescues over $30{,}000$ genuinely stuck
one-leftover-item instances too, confirming \Cref{thm:g2-r1} the same way.
\end{proof}

\subsection{Putting it together}
\label{sec:g2-main}

\begin{theorem}[Conjecture 2 holds at $n=3$]
\label{thm:g2-main}
For every instance of $3$ dichotomous cost functions on a finite item set
$\items$, there is an allocation $\alloc$ and a subsidy $\subsidy \in
\set{0,1}^3$ such that $(\alloc,\subsidy)$ is envy-free, computable in
polynomial time.
\end{theorem}
\begin{proof}
Run the algorithm of \Cref{thm:g2-twyz} to get an EF partial allocation with
at most $2$ leftover items. If none are left, $\subsidy=0$ finishes it. If
one or two items are left, \Cref{thm:g2-r1} or \Cref{thm:g2-r2}
(respectively) produces a full EF allocation with $\subsidy \in
\set{0,1}^3$. Every step is constructive and the extra search in
\Cref{thm:g2-r1,thm:g2-r2} is over a bounded number of choices --- at most
$6$ placements, $6$ permutations, $8$ subsidy vectors --- that does not grow
with $\abs\items$, so the whole procedure runs in polynomial time.
\end{proof}

\paragraph{What this covers, and what it does not.} \Cref{thm:g2-main} makes
no assumption on the cost functions beyond dichotomy: not additive, not
submodular, not related across agents in any way. It is specific to $n=3$:
\Cref{thm:g2-r1,thm:g2-r2} are $3$-agent arguments (the matrices, the
permutation counts, and the subsidy-vector counts are all sized for $3$
agents), even though the starting point, \Cref{thm:g2-twyz}, holds for any
number of agents. Extending the argument to $n \ge 4$ is not attempted here.
